% GPT撰写的独立补充材料，取自其源码预览。
\section{原始执行连接的证明细节}
\label{supp:original-execution}

\noindent\textbf{最早的冲突polka。}
固定一个高度，以及针对非空值$v$在轮次$r$中的提交证书。
从该证书中选取两名诚实签署者组成$K$。
每名签署者均在发出precommit之前锁定$v$。
证书可以在这些签署者已经推进至后续轮次之后组装完成；其被观察到的时点与论证无关。

假设存在更高轮次中针对不同值的polka；若nil允许解锁，也将其计入。
选择存在这类polka的最小轮次$s>r$，并在$s$内选择其不同prevote已经存在的最早事件$\tau$。
其三名投票者中包含一名成员$m\in K$。
由于诚实参与者的轮次递增，$m$在轮次$s$的prevote晚于其在轮次$r$的precommit。
没有解锁依据时，$m$必须对其锁定值发出prevote。
在中间轮次重新锁定$v$仍保持这一约束。
来自低于$s$轮次的冲突依据与$s$的选择矛盾；
来自轮次$s$的依据则必须在$m$发出计入$\tau$的那一票之前已经存在，与$\tau$的选择矛盾。
高于$s$轮次的证据不能先于随后发出的轮次$s$投票。
因此，这样的polka不存在，而更高轮次的冲突提交需要这样的polka。
同轮冲突提交要求某个诚实参与者发出两次precommit。
这一论证依赖所述锁定状态转移及轮次单调规则，而非及时收到提交证据。

\noindent\textbf{配对、验证与执行。}
对候选更新进行归纳：本地构造将正文与其原始分片配对；
远端完成时，根据已检查的分片重建正文。
转移Locked或Valid时，两者共同转移；切换目标时清除旧正文。
历史物化不修改这些活动候选。
针对非空值的precommit在匹配的polka之后产生。
诚实prevote的依据是实际提案验证或更早的锁定；
沿后者回溯时，锁定轮次递减，并最终归于验证。
最终的完整区块身份守卫在Finalize之前固定原始交易、高度、时间及其他应用上下文。
同步与重放提供相同的原始材料。
因此，在这一调用边界内，同高度重试不能通过稳定哈希缓存引入第二份原始正文。

\noindent\textbf{两种等价关系的含义。}
批处理精化比较整批执行后的经济状态与按指定顺序执行各单项步骤后的状态。
该比较通过投影去除命令身份、修订计数器和物化任务。
重放一致性则固定整个原始命令序列及其区块上下文，再由确定性执行重现每个应用哈希。
这是两种不同的关系，二者均不比较以不同顺序安排父交易到达与赔付的历史。正文的延迟中性定理在其自身前提下作这种比较，等同的是最终 CAL 持有人与备付余额，而不是应用哈希或修订元数据。
